\section{Situated Urban and Societal Substrate}
\label{sec:substrate}

GAWorld v2 places persistent agents in a shared substrate whose components
read and write overlapping state.  The city constrains movement, the indoor
spatial tree constrains co-location, local occupancy changes perception,
relationships affect encounters, financial or life events alter later
choices, organizations impose rules, governance proposes and votes on rule
changes, and information exposure shapes what residents believe.  This
coupling supports questions that cross individual and aggregate levels, but
it also means that mechanism provenance and configuration are part of every
experimental treatment.

\subsection{City graph and indoor spatial tree}

The city map parses places, roads, transit lines, interiors, and land-use
attributes into a graph.  Nodes carry categories, coordinates, capacity,
opening hours, popularity, density, and display metadata.  Destination
resolution can match activity terms to place categories, while graph routing
and transport-mode rules estimate distance, time, and fare.  Mode speeds,
fixed delays, road-class multipliers, rush-hour penalties, area price levels,
and weather preferences are explicit parameters.  They make travel costs
inspectable; their presence does not imply transport-model calibration.

v2 extends the city graph with an \emph{indoor spatial tree} that resolves
every indoor activity to a specific room in a building.  The tree shares its
layout vocabulary with the front-end visualizer, so that back-end simulation
and front-end display agree on which room an agent is in.  Eighteen layouts
(residential, apartment, hotel, hospital, clinic, pharmacy, school, library,
shop, cafe, restaurant, office, government, factory, station, gym, park, and
public space) map activity terms to furniture slots (sleep, cook, exercise,
\ldots).  When the spatial-tree plugin is enabled, indoor co-location
restricts a meeting to the same household and the same room; venue capacity
is per-room, computed as a fraction of the venue's capacity that the room
occupies in its layout.  When the plugin is disabled, the v1 behavior is
preserved.

At each tick, the local-physical layer recomputes occupancy from stationary
agents, clears stale counts, and exposes the current location's crowding,
opening status, and weather.  A crowd anomaly is flagged when occupancy is
both high and sharply above the previous tick.  These snapshots can enter
perception and interruption generation, linking collective movement to
individual behavior.  The mechanism is deterministic given agent locations
and parameters, although the locations themselves may follow stochastic or
LLM-mediated decisions.

\subsection{Environment, weather, and crowd anomalies}

The broader environment system supports configured natural, economic,
political, technological, and intraday events.  In rule mode it samples from
configured event distributions using a local random generator.  In LLM mode
a model proposes a daily summary, events, and intraday rules that are
parsed into the same bounded event schema; configured fallback generation
remains available if that call fails.  Severity and anomaly tagging are
rule-based.  v2 separates \emph{natural} events (weather, accident,
epidemic) from \emph{social} events (rumor, market shock, policy reform) and
exposes a \emph{crowd anomaly} channel that the local-physical layer writes
into a single perception field.  Reports must therefore identify the
generator mode rather than treating all environmental histories as
equivalent.

\subsection{Social ties, life events, and offline citizens}

Relationships are typed records rather than undifferentiated neighbor edges.
Role templates specify channels, obligation, protection, and daily decay;
interaction signals update closeness and contact time.  The implementation
can bootstrap off-screen contacts from an LLM-generated backstory, with a
deterministic heuristic fallback.  v2 adds an \emph{offline citizen}
subsystem that classifies each ghost contact by relationship role and
admission channel (call, message, visit, none), produces stateful obligations
when a contact is long-neglected, and exposes a visit channel that the
\emph{travel} plugin \citep{gaworldtravel} actually exercises.  The combined
offline + travel stack turns a previously read-only field (the obligation
counter) into a behaviorally consequential signal.

Maintenance is rule-based.  Uncontacted ties decay at role-specific rates,
long-neglected obligations can rise, and an enforcement pass assigns nested
Dunbar-style tiers while pruning beyond a configured outer cap and protecting
selected kin roles.  These tiers are an operational capacity constraint, not
evidence that the simulated network reproduces measured human networks.

Life events are normalized, persisted, scheduled by simulated day and time,
and drained once when due.  Events can target agents and carry bounded state
effects and tags.  v2 makes life-event state effects typed and reversible:
each effect has an effect-id, a magnitude, a half-life, and a reversal hook,
so a cancelled event can be undone at the day boundary without leaving
free-floating state.  This is a controlled exogenous-event surface whose
timing and direct updates are inspectable; the realism of event frequencies,
dependencies, and consequences must be established by a separate study.

\subsection{Economy, credit, and endogenous inflation}

The finance module maintains checking, savings, investment, and housing-fund
accounts; monthly budgets; income and expense records; liabilities and
credit; sector pools; government flows; and macro state.  Configured rules
calculate social-insurance deductions and progressive income tax, route wages
and consumption, enforce available-cash and credit constraints, transfer
savings, and settle payments.  A dedicated seeded random generator drives
expenditure variation, market returns, shocks, and other probabilistic
transitions so that unrelated random draws need not perturb the economy
sequence.

v2 closes a long-standing gap: the v1 code reported an inflation rate and a
price index that no other component read.  In v2 the price level multiplies
real expenditure so that wages that do not follow it produce a measurable
erosion of real consumption.  Inflation itself is endogenized through a
\emph{demand-pull} rule: $\pi = \pi^{*} + \kappa \cdot (C/C_0 - 1)$, where
$C$ is a 30-day moving average of real per-capita consumption, $C_0$ is the
first 30 days as a baseline, and $\pi$ is clamped to a configured interval
\citep{gaworldinflation}.  The wage index, when enabled, ties a fraction of
new wages to the price level.  This is the first run at making
``cost-of-living erosion'' actually measurable in a GAWorld run, and the
artifact ledger records whether the wage index was active so that the
directional claim is not silently over-interpreted.

Shared sector and government pools make counterparties visible instead of
creating every monetary flow without an account.  Agent-to-agent payments
and some rent or labor flows can be routed to modeled recipients;
investment returns combine a common market component with idiosyncratic
variation.  Ledger and conservation checks can reveal implementation defects
such as unbalanced internal transfers.  They cannot establish that the tax
schedule, consumption response, labor process, credit behavior, or macro
dynamics is an empirically valid representation of a real economy.
Accounting consistency is therefore a software invariant, not a
behavioral-validation result.

\subsection{Organizations, governance, and rule changes}

v2 adds a persistent-organization subsystem.  Each organization is a
typed object (community, business, club) with members, role templates, a
chosen governance mode (\emph{off}, \emph{leader}, \emph{member-vote}), and
a ruleset that the plugin applies on member decisions.  Two rule-change
pipelines exist: a \emph{leader} pipeline in which a designated member
decides, and a \emph{member-vote} pipeline in which a proposal is opened,
voted on, and (if it passes) enqueued for application at the next day
boundary \citep{gaworldgovernance}.  Generations are isolated: a proposal
opened under a given member snapshot, governance version, and ruleset cannot
be re-opened after the world is rebuilt with a different member snapshot.
This isolation is enforced in a single SQLite transaction.

The governance subsystem is opt-in (\texttt{governance.enabled}); when it is
off, the v1 organization behavior is preserved.  Vote tallies and rule
versions are recorded in a separate \texttt{governance.json} per generation,
not interleaved with the per-step metrics CSVs, so that governance results
do not contaminate aggregate economic measurements.

\subsection{Information sources and Moltbook}

v2 unifies information exposure under an \emph{info-sources} surface
\citep{gaworldrag}.  Each resident has a configured diet (a list of
channels and a sampling rate) and a per-tick search budget.  Channels cover
local news, social posts, professional feeds, and (when enabled) web
search.  The Moltbook integration gives residents a social network identity
(\texttt{gaworld/moltbook}) where they post, reply, and react; the agent's
actions in Moltbook are recorded as separate \texttt{log.jsonl} entries,
distinct from simulation traces, and are summarized back into the perception
field at the next tick.  This separation lets a researcher inspect the
social-network behavior of a resident without conflating it with the
simulation's per-step state history.

\subsection{Coupling and ownership boundaries}

\Cref{tab:component-coupling} makes the principal shared-state paths
explicit.  ``Implemented'' refers to an inspectable current mechanism;
``partial boundary'' indicates that the mechanism is connected but retains
direct calls, shared dictionaries, or specialized writers rather than a
fully isolated plugin interface.

\begin{table}[t]
\centering
\small
\setlength{\tabcolsep}{3.5pt}
\caption{Coupling among agent and societal components.  Evidence status
describes implementation visibility, not empirical validity.}
\label{tab:component-coupling}
\begin{tabular}{@{}p{0.13\linewidth}p{0.16\linewidth}p{0.16\linewidth}p{0.27\linewidth}p{0.16\linewidth}@{}}
\toprule
Subsystem & Reads & Writes & Agent-visible consequence & Evidence status \\
\midrule
Memory and cognition & Episodes, state, context, retrieved entries & Memories, intentions, habits, relationships & Recalled context, priorities, reflection, later action bias & Implemented; human realism unvalidated \\
Interests and skills & Profile, episodes, social growth context & Growth profile, private skills & Practice focus and procedural prompt context & Implemented; partial plugin boundary \\
World and local physics & City graph, time, weather, locations, rooms & Routes, travel state, occupancy snapshots, room occupancy & Feasible destinations, travel costs, crowding, closure cues, room-level co-location & Implemented; direct integration \\
Environment & Configuration, history, optional LLM output & Events, weather and macro context & Interruptions and perceived exogenous conditions & Implemented; generator-mode dependent \\
Social network & Profiles, contact history, interactions & Typed ties, closeness, obligation, tiers & Encounter probability and relationship context & Implemented; mechanism validity unassessed \\
Life events & Persisted event queue, clock & Consumption flag and configured state effects (typed, reversible) & Scheduled exogenous changes to targeted agents & Implemented; committed plugin-integrated event producer \\
Economy & Activity, balances, calendar, sectors, macro state, price level & Accounts, ledger, budgets, pools, agent state, end-of-month price index & Income, spending constraints, taxes, credit, returns, payments, real-consumption erosion & Implemented; primarily hook-based \\
Work & Capability cache, activity, growth, market jobs & Queue results, artifacts, settlement and state effects & Produced files, earnings, experience records & Implemented; artifact quality unvalidated \\
Indoor rooms and capacity & City map, location, activity & Per-room occupancy, capacity enforcement, room-intent binding & Room-resolved perception, capacity-based rejection & Implemented; opt-in \\
Personality & Profile, optional LLM overlay & Personality proxies & Conditional interruption probability, growth-phase weighting & Implemented; anchor + overlay \\
Organizations and governance & Membership, role, ruleset, proposal & Updated ruleset, accepted proposal, vote tallies & Rule-conformant behavior, post-vote transitions & Implemented; opt-in \\
Offline citizens and travel & Off-screen contact roles, duty obligations & Visit channel, trip itinerary, location rewrite & Mobile-agent perception, family duty fulfillment & Implemented; opt-in \\
Information sources & Channel list, search budget, agent-derived query & Retrieved entries, Moltbook posts, perception field & Exposure heterogeneity, search-induced awareness & Implemented; channel-dependent \\
\bottomrule
\end{tabular}
\end{table}

The table exposes the current modularity limit.  The 14 built-in plugins
have plugin entry points; economy, local physical perception, spatial
preferences, work, dynamic behavior, indoor rooms, organizations,
governance, personality, offline citizens, and information sources have
built-in adapters.  Memory, environment evolution, and social-network
operations still exchange mutable agent dictionaries or runtime objects
directly.  This arrangement is usable and inspectable, yet replacing one
mechanism can change prompts, state schemas, random draws, and specialized
artifacts elsewhere.  Controlled comparisons should therefore freeze the
complete coupled configuration rather than naming only the focal
intervention.
